\documentclass[11pt,a4paper]{article}
\usepackage[margin=25mm]{geometry}
\usepackage{amsmath,amssymb,amsthm,hyperref}
\newtheorem{lemma}{Lemma}
\newcommand{\gen}[1]{\langle#1\rangle}
\title{Fixed prefixes and support deficits at a unique orbit crossing}
\author{Human Compression project: L044 working continuation}
\date{}
\begin{document}\maketitle
Let $G=\gen{\rho_0,\ldots,\rho_{r-1}}$ be a faithful transitive
string action by involutions. Suppose colour $f$ has a unique edge
$\{a,b\}$ whose endpoints are in different orbits of
$G_f=\gen{\rho_j:j\ne f}$. The single-crossing lemma in
\emph{Paired orbit cuts in string actions} gives exactly two
$G_f$-orbits $A,B$, containing $a,b$, and says that every
$\rho_j$ with $|j-f|>1$ fixes both $a,b$. Write
$G_{>l}=\gen{\rho_j:j>l}$ and define $G_{<l}$ dually.

\begin{lemma}[A fixed prefix to an active colour]\label{prefix}
Suppose $l>f$ and $\rho_l$ acts nontrivially on $A$. There is a
path in the full permutation graph from $a$ to a point moved by
$\rho_l$, using only colours less than $l$ and different from $f$.
Every vertex of this path is fixed by $G_{>l}$.
The dual assertion for $l<f$ uses only colours greater than $l$
and different from $f$, and its vertices are fixed by $G_{<l}$.
The same assertions hold with $b,B$ in place of $a,A$.
\end{lemma}
\begin{proof}
Put $H=\gen{\rho_j:j<l,\ j\ne f}$ and $R=G_{>l}$.
Every generator of $R$ fixes $a$, because $j>l>f$ implies
$j-f\ge2$. Every generator of $R$ also commutes with every
generator of $H$. Consequently $R$ fixes the whole orbit $a^H$
pointwise. If $\rho_l$ fixed that orbit pointwise as well, then
the orbit would be invariant under all the generators of $G_f$.
It would therefore equal $A=a^{G_f}$, contradicting the assumed
nontrivial action of $\rho_l$ on $A$. Some $H$-word thus takes
$a$ to a point moved by $\rho_l$. Its successive vertices lie
in $a^H$ and are fixed by $R$. Erasing loops gives the required
simple path. Reversing the index order proves the dual claim;
the argument at $b$ is identical.
\end{proof}

\begin{lemma}[Consecutive-colour lower bound]\label{length}
Under the hypotheses of Lemma~\ref{prefix} with $l>f$, every
path from $a$ to a point moved by $\rho_l$ that avoids colour $f$
contains each of $f+1,\ldots,l-1$. In particular appending the
final $l$-edge to a shortest prefix in Lemma~\ref{prefix} gives
a simple path with at least $l-f$ edges. Its prefix, excluding
that final edge, is pointwise fixed by $G_{>l}$.
\end{lemma}
\begin{proof}
If a colour $k$ with $f<k<l$ were missing, factor the path word
into a word in generators of index less than $k$ and a word in
generators of index greater than $k$. The two sets commute.
Every generator of index greater than $k$ fixes $a$, because
its index is at least $f+2$. Thus the endpoint is reached from
$a$ by the lower-index word alone. The generator $\rho_l$
commutes with that word and fixes $a$, so fixes the endpoint,
a contradiction. This proves the colour and length assertions.
Choose a shortest path, among the permitted prefixes of
Lemma~\ref{prefix}, to any point moved by $\rho_l$. Every
earlier vertex is then fixed by $\rho_l$. The other endpoint
of the final $l$-edge is also moved by $\rho_l$, so cannot be
an earlier prefix vertex. The extended path is simple.
\end{proof}

\begin{lemma}[The endpoint case of source Proposition 5.19]\label{endpoint}
Suppose $f=0$, every generator is even, $r\ge4$, and
$\rho_{r-1}$ acts nontrivially on both $G_0$-orbits. Then
$n\ge2r+1$, equivalently $r\le\lfloor(n-1)/2\rfloor$.
\end{lemma}
\begin{proof}
Apply Lemma~\ref{length} in each orbit with $l=r-1$. The two
disjoint simple paths each have at least $r-1$ edges, so together
they have at least $2r$ vertices. It remains to exclude $n=2r$.

At equality the paths exhaust their respective orbits. Each has
exactly one edge of each colour $1,\ldots,r-1$, in that order:
the no-gap assertion applied to each initial segment forces the
first occurrence of every colour to follow all smaller positive
colours. Write the paths
\[
 a_0=a\mathrel{-_1}a_1\mathrel{-_2}\cdots
 \mathrel{-_{r-1}}a_{r-1},\qquad
 b_0=b\mathrel{-_1}b_1\mathrel{-_2}\cdots
 \mathrel{-_{r-1}}b_{r-1}.
\]
Because $a_0b_0$ is the unique cross-orbit $0$-edge,
$\rho_0$ preserves each tail $A'=\{a_1,\ldots,a_{r-1}\}$ and
$B'=\{b_1,\ldots,b_{r-1}\}$. The group
$H=\gen{\rho_2,\ldots,\rho_{r-1}}$ fixes $a_0,b_0$ and
preserves both tails; its action on each tail is transitive,
since the displayed edges form a path through that tail.

The subgroup $K=\gen{\rho_3,\ldots,\rho_{r-1}}$ fixes $a_1$:
each of its generators fixes $a_0$ and commutes with $\rho_1$.
Every other vertex of $A'$ is moved by one of the displayed
edges of colour at least three. This uses $r\ge4$, and shows
that $a_1$ is the unique $K$-fixed point in $A'$. The restriction
of $\rho_0$ centralises $H$, hence $K$, so must fix $a_1$.
Centralising the transitive $H$-action then forces it to fix
every point of $A'$. The identical argument fixes all of $B'$.
Therefore $\rho_0=(a_0\,b_0)$ is a single transposition,
contrary to evenness. This excludes equality.
\end{proof}

For the main L044 range $n\ge12$, ranks $r\le3$ already satisfy
the required bound, so Lemma~\ref{endpoint} covers the whole
endpoint/top-colour case. The argument does not assume full IP,
does not use the unresolved square-rigidity assertion, and does
not assert that all other source branches are complete.


\begin{lemma}[The side of a shared-colour tail]\label{side}
Suppose $f=0$, the noncommutation diagram of the given marking is
connected, and $h<r-1$ is the largest index whose generator acts
nontrivially on both $A$ and $B$. All generators of index greater
than $h$ act nontrivially on just one and the same orbit, say $B$.
In particular $G_{>h}$ fixes $A$ pointwise.
\end{lemma}
\begin{proof}
Each generator of index greater than $h$ is nonidentity and acts on
at most one of the two $G_0$-orbits. It therefore acts on exactly
one. If consecutive generators in this tail acted on different
orbits, their supports would be disjoint, so they would commute.
That would disconnect the noncommutation diagram: all
nonconsecutive generators already commute by the string condition.
Thus the entire tail acts on the same orbit.
\end{proof}

\begin{lemma}[A support deficit without the source orbit reversal]
\label{deficit}
Under the hypotheses of Lemma~\ref{side}, suppose $h\ge1$.
Put $X=\operatorname{supp}(G_{>h})$. Then
\[
 X\subseteq B,\qquad |\operatorname{Fix}(G_{>h})|\ge2h+1,
 \qquad h\le\frac{n-|X|-1}{2}.
\]
More precisely, the fixed set contains the disjoint union of
an $A$-path from $a$ including its final $h$-edge, with at least
$h+1$ vertices, and the prefix of a $B$-path from $b$ to an
$h$-edge, with at least $h$ vertices.
If equality holds, these paths have exactly the indicated sizes
and exhaust the fixed set; $A$ is the first path's vertex set.
The paths are ordered by the colours $1,\ldots,h$.
This last assertion is a description of the specified paths and
fixed set, not a classification of all full-graph edges.
\end{lemma}
\begin{proof}
The $A$-path supplied by Lemma~\ref{length}, including its last
$h$-edge, is fixed by the whole tail because that tail fixes all
of $A$. It has at least $h+1$ vertices. The $B$-prefix is fixed
by Lemma~\ref{prefix}, and has at least $h$ vertices. These sets
are disjoint, proving the bound. At equality the two lower bounds
are exact and account for every fixed vertex. Since every point
of $A$ is fixed, no other point of $A$ can remain. The first
occurrence/no-gap argument used in Lemma~\ref{endpoint} orders
the path colours. There is no claim that the final $B$-endpoint
is fixed: its omission is precisely why the count is $2h+1$.
\end{proof}

For the L044 application, a nontrivial faithful involutory marking
of the simple group $A_n$ ($n\ge5$) has connected noncommutation
diagram. Otherwise two nonempty sets of nonidentity generators
would generate commuting normal factors of $A_n$; simplicity and
its trivial centre exclude that. Thus Lemma~\ref{side}'s explicit
hypothesis is available. No C-group intersection argument is
needed for these support bounds.

The main inequality and support placement in
Lemma~\ref{deficit} give the useful part of source
Proposition~5.20. Its literal statement places $X$ in the orbit
\emph{fixed} by $\rho_{r-1}$, while its proof defines $X$ as the
support of the tail containing $\rho_{r-1}$. The actual side is
$B$, on which that generator acts. The source's displayed
whole permutation-graph assertion has not been replaced by an
unproved claim: only the exact path/fixed-set consequence above
is established here. Its later use must preserve that boundary.


\begin{lemma}[Exact exceptional end under an alternating tail]
\label{alternatingtail}
In Lemma~\ref{deficit}, suppose the marking satisfies IP, all
its generators are even, $G_{>h}$ is the full natural alternating
group on $X$, and $|\operatorname{Fix}(G_{>h})|=2h+1$.
Then $h=2$, and $G_{<3}$ is exactly the six-point marking in
source Table~2(2), fixing every other point.
\end{lemma}
\begin{proof}
The nontrivial alternating group generated by the involutions of
the tail acts on at least five points: $A_3$ is odd-order and the
involutions of $A_4$ generate only its Klein four subgroup.
The centraliser of the natural $A_X$ in $\operatorname{Sym}(X)$
is trivial. For example its point stabiliser fixes only that point,
so an element of the centraliser fixes one point and then all points
by transitivity. Since $\rho_0$ commutes with the tail, it fixes $X$
pointwise.

Write the equality paths as in Lemma~\ref{endpoint}, stopping at
index $h$. The endpoint $b_h$ lies in $X$: the rest of the two paths
already exhaust the fixed set. Thus $\rho_0$ fixes $b_h$.
Commutation with $\rho_h,\rho_{h-1},\ldots,\rho_2$ propagates
this fixation to $b_{h-1},\ldots,b_1$. Hence its only possible
remaining action, besides $(a_0\ b_0)$, is on
$A'=\{a_1,\ldots,a_h\}$.

If $h=1$, that set is a singleton, contradicting evenness of
$\rho_0$. If $h\ge3$, the group
$\gen{\rho_2,\ldots,\rho_h}$ is transitive on $A'$, and its
subgroup $\gen{\rho_3,\ldots,\rho_h}$ has exactly one fixed
point there, namely $a_1$. The proof is the same displayed-path
argument as in Lemma~\ref{endpoint}. Centralising both groups
forces $\rho_0$ to fix $A'$ pointwise, again a contradiction.
Therefore $h=2$ and evenness gives
\[
 \rho_0=(a_0\ b_0)(a_1\ a_2).
\]
Now $\rho_1$ centralises the whole tail, so fixes $X$ pointwise.
On the three-point orbit $A$ and the two fixed points of $B$ it is
therefore exactly
\[
 \rho_1=(a_0\ a_1)(b_0\ b_1).
\]
The generator $\rho_2$ fixes $a_0,b_0$ and has the form
\[
 \rho_2=(a_1\ a_2)(b_1\ x)\delta,\qquad x=b_2\in X,
\]
where $\delta$ is an even involution supported on
$X\setminus\{x\}$. On the six displayed points the three core
permutations generate the natural degree-six action of $A_5$;
the complete small certificate \texttt{small\_tail\_model.py}
identifies them, generator by generator, with three permutations
in $A_5\le S_5$. This is the source's Table~2(2), with its vertices
read left to right as $a_2,a_1,a_0,b_0,b_1,x$.

If $\delta\ne1$, the prefix $G_{<3}$ is a subgroup of the direct
product of this core $A_5$ and $\gen{\delta}$, projecting onto
both factors. If it did not contain $\delta$, it would be the
graph of a homomorphism $A_5\to C_2$ taking the third core
involution to the nonidentity element. Perfection of $A_5$
excludes such a homomorphism. Thus $\delta\in G_{<3}$; but
$\delta\in A_X=G_{>2}$ as well, contrary to IP for the disjoint
index sets $\{0,1,2\}$ and $\{3,\ldots,r-1\}$. Consequently
$\delta=1$, proving the exact marked conclusion.
\end{proof}

This lemma supplies the endpoint alternating-tail equality
obstruction used in the source's Proposition~5.26. It does not
replace the other prefix cases, the full Table~2 census, or the
remaining single-crossing induction.


\begin{lemma}[A transitive tail with an earlier commuting involution]
\label{transitivetail}
Assume every maximal omission subgroup $G_k$ is intransitive.
Let $H=G_{>j}$ be transitive on its support $X$, and suppose
$\rho_l$, for some $l<j$, acts nontrivially on $X$. Then
\[
 r-1-j\le |X|/2-1.
\]
No IP is required for this inequality. It uses the retained Whiston
bound for an independent family in a symmetric group.
\end{lemma}
\begin{proof}
The restriction $z=\rho_l|_X$ commutes with $H$. An element
centralising a transitive group and fixing one point fixes all
points. Hence the nontrivial involution $z$ has no fixed point,
and its pairs give an $H$-invariant block system with
$m=|X|/2$ blocks.

Choose a minimal subset $L$ of the displayed $H$ generators whose
images generate the whole block action. The group $\gen{L}$ is
transitive on these blocks and has either one or two orbits on
$X$. If there are two, $z$ interchanges them (each meets every
pair in one point). Since $H$ itself is transitive, there is a
generator $\rho_k$ of $H$ outside $L$. Then $G_k$ contains both
$\gen{L}$ and $\rho_l$, and is transitive on $X$. If instead
$\gen{L}$ is already transitive, $G_k$ is transitive on $X$ for
any $H$ generator $\rho_k$ outside $L$.

In either case an omitted $\rho_k$ is supported entirely in the
single $G_k$-orbit containing $X$. Adjoining it cannot join any
$G_k$-orbits. Since $\gen{G_k,\rho_k}=G$ is transitive, $G_k$
would already be transitive, a contradiction. Thus every $H$
generator belongs to $L$. Their block images form an independent
family in $S_m$, so their number is at most $m-1$ by Whiston.
\end{proof}

Lemma~\ref{transitivetail} supplies the rank inequality of source
Proposition~5.21. Its displayed equality graph classification is
a separate obligation; it is not inferred from Whiston's cardinality
bound alone.

\paragraph{Exact relation to the source.}
The no-gap idea comes from Cameron--Fernandes--Leemans--Mixer,
\emph{Highest rank of a polytope for $A_n$},
arXiv:1605.09173v1. The relevant source assertion is
Proposition~5.18, on physical pages~29--30.
The statement above specifies the initial orbit, the active target
colour, the side of the string which fixes the path, and the
exclusion of its final edge. The source's final sentence instead
prints $G_{>l}$ when $l<f$ and $G_{<l}$ when $l>f$ and does not
make the prefix boundary explicit. We do not identify that literal
sentence with Lemma~\ref{prefix}. No bound for the whole original,
new conjecture solution or formal verification is asserted here.
\end{document}
